% Physical unit 017; accepted Session 02.
\GordanUnitStart{017}{printed page 12}
\GordanTopFolio{12}
\[
\begin{aligned}
&=\int_{\psi_1}^{\psi_2}d\psi\,
{\begin{gathered}
\delta l'\delta l'\cos^2l(1-e^2\cos^2l)
+2\delta l\delta l'\sin l\cos l\,l'\bigl(e^2l'^2[1-e^2\cos^2l]+[1+e^2\cos^2l]\bigr)\\
+\delta l^2\bigl(e^2l'^4[\cos2l-e^2\cos^4l]+l'^2[2e^2\cos^4l-2\cos2l]\bigr)
\end{gathered}
\over \{(1-e^2\cos^2l)l'^2+\cos^2l\}^{1\frac12}}\\
&=\int_{\psi_1}^{\psi_2}d\psi\{(11)\delta l'\delta l'+2(01)\delta l\delta l'+(00)\delta l\delta l\}>0
\end{aligned}
\]
where, for this integral, $\delta l=0$ at the limits.

\begin{center}\S\ 4.\end{center}

The elements of this integral, for values assigned to the variation $\delta l$, can become either positive or negative. Whether the whole integral is positive or negative is determined by this, whether the sum of the positive elements or of the negative elements is the greater. Since, however, the question is how far our integral is positive for values of integration of whatever kind, the whole question must be referred to this: how far an integral is found which, since it is itself $0$, when added to that integral, makes all its elements positive in force.

\begin{center}\S\ 5.\end{center}

An integral of this kind is the following:
\[
 \int_{\psi_1}^{\psi_2}d\psi\{\delta l\delta l N\}'=
 \int_{\psi_1}^{\psi_2}d\psi\{2N\delta l\delta l'+N'\delta l\delta l\}.
\]
This is so in so far as the function $N$ can be defined in such a way that the sum
\[
 (11)\delta l'\delta l'+2[(10)+N]\delta l\delta l'+[(00)+N']\delta l\delta l
\]
throughout the whole interval, with $N$ never becoming infinite, becomes a complete square.

For under this condition
\[
 \int_{\psi_1}^{\psi_2}d\psi\{2N\delta l\delta l'+N'\delta l\delta l\}=N_2\delta l_2\delta l_2-N_1\delta l_1\delta l_1=0
\]
since the variations vanish at the limits.

\begin{center}\S\ 6.\end{center}

The sum
\[
 (11)\delta l'\delta l'+2[(10)+N]\delta l\delta l'+[(00)+N']\delta l\delta l
\]
then is a square when
\[
 (11)[(00)+N']=[(01)+N]^2,
\]
or, integrating,
\[
 -N=(01)+(11){\alpha {\partial l'\over\partial\psi_1}+\beta {\partial l'\over\partial l_0}
 \over \alpha {\partial l\over\partial\psi}+\beta {\partial l\over\partial l_0}}.
\]
In this equation $\alpha$ and $\beta$ are arbitrary constants.



% ---- GDL03 continuation, source PDF pages 018--027 ----
